\section{Clasificadores supervisados}
\label{sec:met_supervisados}
% Extensión estimada: 4--4.5 páginas, con una figura y tres tablas, en tres subsecciones.
% Fuentes: src/models/classifiers.py; src/data/transforms.py; src/training/recipes.py, optim_groups.py,
% schedulers.py, losses.py; src/evaluation/tta.py; scripts/train_classifier.py; scripts/toko/submit_sweep.py
% y submit.py (valores efectivos); scripts/sweep_utils.py y configs/sweeps/phase1_modern_sweep_grid.yaml
% (grilla); logs de entrenamiento (parámetros); docs/resul/split_lesion_resul_20260926.md (ganadores, semillas,
% réplica del desbalance); docs/resul/densenet_vit_moderna_resul_20261017.md (DenseNet121 y ViT-B/16 con la
% configuración de su familia, tres semillas); scripts/lesion_split_report.py, classifier_ensemble.py,
% melanoma_detection.py. No se usan como fuente los configs/*.yaml de clasificación, porque train_classifier.py no los lee.
% Revisión del autor del 18/10/2026: solo la configuración final; sin comparar recetas ni describir la búsqueda
% anterior de DenseNet121 y ViT-B/16, cuyas corridas quedan fuera del capítulo.
% Revisión del tutor del 21/10/2026: estructura común con 3.4 (propósito → modelos seleccionados → configuración →
% procedimiento experimental); el motivo de la elección antes de las características; figura de uso; los criterios
% de selección se definen en 3.7 y acá solo se nombran.

Los clasificadores supervisados cumplen dos funciones en el trabajo. Por un lado, resuelven la clasificación (T1) y la
detección de melanoma (T2) con modelos entrenados solo con imágenes, y son la línea base
(Sección~\ref{sec:marco_swin}) con la que se mide si un \gls{vlm} clasifica igual de bien. Por otro lado, dentro del
\emph{pipeline}, son los que fijan el diagnóstico que recibe el generador. Para cumplir ambas funciones hacen falta
modelos que aprendan de unas 6\,400 imágenes y que puedan entrenarse en una GPU de consumo.

\subsection{Modelos seleccionados}
\label{sec:met_supervisados_arquitecturas}

Las dos familias de arquitecturas de visión, las convolucionales y las basadas en transformers, parten de supuestos
distintos sobre la imagen, y ninguna se impone de antemano (Sección~\ref{sec:marco_vision}). A su vez, los trabajos de
clasificación dermatológica revisados evalúan casi solo redes convolucionales (Sección~\ref{sec:estado_arte_diagnostico}).
Por ello, el trabajo compara las dos familias con el mismo procedimiento, y de cada una toma dos de los modelos
presentados en el Capítulo~\ref{cp:marco}: DenseNet121 y ConvNeXt-Tiny entre las convolucionales, y ViT-B/16 y
Swin-Tiny entre las basadas en transformers. Las variantes Tiny, con unos 27.5 millones de parámetros, tienen un tercio
de los de ViT-B/16 (Tabla~\ref{tab:met_modelos_sup}).

Los cuatro modelos se usan de la misma manera, que muestra la Figura~\ref{fig:met_clasificador}. Cada uno parte de
pesos preentrenados en \emph{ImageNet}-1K, y su última capa se reemplaza por una cabeza de siete salidas, precedida
por un \emph{dropout} de 0.3. La softmax convierte esas salidas en una probabilidad por clase, de la que se obtienen
las dos respuestas del clasificador: la clase de mayor probabilidad, para T1, y la probabilidad de melanoma, que se
compara con un umbral para T2.

% Figura: uso de un clasificador supervisado en este trabajo, fig:met_clasificador. Usada en sec:met_supervisados (3.3).
% No redibuja la arquitectura (Cap. 2): muestra entrada, red preentrenada, cabeza nueva, salidas para T1 y T2 y qué se
% entrena en cada etapa. Revisión del tutor del 21/10/2026 (figuras de los modelos). Ancho: unos 148 mm.

\begin{figure}[H]
    \centering
    \hyphenpenalty=10000 \exhyphenpenalty=10000
    \begin{tikzpicture}[x=1mm, y=1mm]

        \node[figentrada, text width=15mm] (img) at (9, 0) {imagen\\{\footnotesize $224 \times 224$}};
        \node[figtarea, text width=32mm] (red) at (42, 0)
            {red preentrenada\\{\footnotesize DenseNet121, ConvNeXt-Tiny, ViT-B/16 o Swin-Tiny}};
        \node[figtarea, text width=21mm] (cab) at (79, 0) {cabeza nueva\\{\footnotesize \emph{dropout} y lineal de 7 salidas}};
        \node[fignodo, text width=17mm, fill=figrelleno] (sm) at (105, 0) {softmax\\{\footnotesize $p_c$ por clase}};
        \node[figsalida, text width=21mm] (t1) at (135, 9) {T1: clase de mayor $p_c$};
        \node[figsalida, text width=21mm] (t2) at (135, -9) {T2: $p_{\text{MEL}}$ frente a un umbral};

        \draw[figflecha] (img) -- (red);
        \draw[figflecha] (red) -- (cab);
        \draw[figflecha] (cab) -- (sm);
        \draw[figflecha] (sm.east) -- ++(4, 0) |- (t1.west);
        \draw[figflecha] (sm.east) -- ++(4, 0) |- (t2.west);

        \node[figetiqueta, align=center, text=fignaranja!70!black, anchor=north] at ([yshift=-2mm]red.south)
            {etapa 1: congelada\\etapa 2: se ajusta toda};
        \node[figetiqueta, align=center, text=fignaranja!70!black, anchor=north] at ([yshift=-2mm]cab.south)
            {se entrena en\\las dos etapas};
        \node[figetiqueta, anchor=south] at ([yshift=2mm]red.north) {pesos de \emph{ImageNet}-1K};

    \end{tikzpicture}
    \caption{Uso de los clasificadores supervisados: la red preentrenada recibe una cabeza nueva, y la probabilidad por
    clase resuelve T1 y T2.}
    \label{fig:met_clasificador}
\end{figure}


\begin{table}[H]
    \centering
    \caption{Los cuatro clasificadores: familia, preentrenamiento, cabeza, parámetros con cabeza de siete clases y tamaño de lote.}
    \label{tab:met_modelos_sup}
    \small
    \begin{tabularx}{\textwidth}{@{}>{\raggedright\arraybackslash}p{2.1cm}>{\raggedright\arraybackslash}X>{\raggedright\arraybackslash}p{1.9cm}>{\raggedright\arraybackslash}p{2.6cm}>{\raggedleft\arraybackslash}p{1.85cm}>{\raggedleft\arraybackslash}p{0.8cm}@{}}
        \toprule
        \textbf{Modelo} & \textbf{Familia} & \textbf{Pesos iniciales} & \textbf{Cabeza} & \textbf{Parámetros} & \textbf{Lote} \\
        \midrule
        DenseNet121 & Red convolucional, bloques densos & \emph{ImageNet}-1K & \emph{Dropout} 0.3 y lineal de 7 salidas & 6\,961\,031 & 32 \\
        \filasep
        ConvNeXt-Tiny & Red convolucional moderna & \emph{ImageNet}-1K & \emph{Dropout} 0.3 y lineal de 7 salidas & 27\,825\,511 & 32 \\
        \filasep
        ViT-B/16 & Transformer visual, parches de $16 \times 16$ & \emph{ImageNet}-1K & \emph{Dropout} 0.3 y lineal de 7 salidas & 85\,804\,039 & 16 \\
        \filasep
        Swin-Tiny & Transformer jerárquico, con ventanas & \emph{ImageNet}-1K & \emph{Dropout} 0.3 y lineal de 7 salidas & 27\,524\,737 & 16 \\
        \bottomrule
    \end{tabularx}
\end{table}

\subsection{Configuración del entrenamiento}
\label{sec:met_supervisados_receta}

La configuración parte de la entrada. Las imágenes de \emph{HAM10000} se reducen una sola vez a $224 \times 224$
píxeles, sin conservar su proporción, y todos los modelos de este capítulo parten de esa versión. Los clasificadores,
además, las reciben normalizadas con la media y el desvío de \emph{ImageNet}. Como el entrenamiento cuenta con 6\,447
imágenes (Sección~\ref{sec:met_particiones}), el riesgo principal es el sobreajuste
(Sección~\ref{sec:marco_regularizacion}). Para reducirlo se aplican aumentos de datos que cambian la imagen sin
cambiar su clase, ya que una misma lesión puede fotografiarse con otro encuadre, otra orientación u otra iluminación.
En concreto, se toma un recorte aleatorio de entre el 70\,\% y el 100\,\% del área, se voltea la imagen horizontalmente
con probabilidad 0.5, se aplican dos transformaciones de RandAugment~\cite{cubuk2020randaugment} con una magnitud $M$
que se elige por búsqueda, y se borra un rectángulo al azar con probabilidad 0.25~\cite{zhong2020random}.

En validación y en prueba, en cambio, no hay aumento. La imagen se redimensiona a $256 \times 256$ y se recorta el
centro de $224 \times 224$. Las probabilidades finales de ambas particiones se obtienen con \gls{tta}
(Sección~\ref{sec:marco_regularizacion}), que promedia las de la imagen y su reflejo horizontal, mientras que la
validación de cada época usa una sola vista.

Con las entradas fijadas, cada modelo se adapta a \emph{HAM10000} en dos etapas, que corresponden a las dos opciones
de la Sección~\ref{sec:marco_transferencia}: primero la extracción de características y después el ajuste fino. En la
primera etapa, con el resto de la red congelado, se entrena solo la cabeza nueva durante una época, de modo que el
ajuste fino no parta de una cabeza inicializada al azar. En la segunda se descongela toda la red, y cada capa recibe
0.85 veces la tasa de aprendizaje de la capa siguiente (Ecuación~\ref{eq:marco_llrd}). Esta etapa dura hasta 39
épocas, con tres de calentamiento lineal y decaimiento coseno de la tasa~\cite{goyal2017accurate,loshchilov2017sgdr}.

Los demás elementos del entrenamiento son comunes a los cuatro modelos. El optimizador es
AdamW~\cite{loshchilov2019decoupled}, y la función de pérdida es la entropía cruzada con suavizado de etiquetas de
0.1~\cite{szegedy2016rethinking,muller2019does}. Finalmente, la detención temprana corta el entrenamiento si el F1
ponderado de validación no mejora en cinco épocas, y se conserva la mejor época; la
Sección~\ref{sec:met_seleccion} justifica ese criterio. La Tabla~\ref{tab:met_receta} reúne estos valores.

\begin{table}[H]
    \centering
    \caption{Parámetros de entrenamiento comunes a los cuatro clasificadores.}
    \label{tab:met_receta}
    \small
    \begin{tabularx}{\textwidth}{@{}>{\raggedright\arraybackslash}p{4.2cm}>{\raggedright\arraybackslash}X@{}}
        \toprule
        \textbf{Parámetro} & \textbf{Valor} \\
        \midrule
        Optimizador & AdamW \\
        \filasep
        Épocas & Hasta 40: una de cabeza y hasta 39 de red completa \\
        \filasep
        Decaimiento de la tasa por capa & 0.85 \\
        \filasep
        Planificador & Calentamiento lineal de tres épocas desde 0.1 de la tasa; luego coseno \\
        \filasep
        Función de pérdida & Entropía cruzada, con suavizado de etiquetas de 0.1 \\
        \filasep
        Detención temprana & Paciencia de cinco épocas sobre el F1 ponderado de validación \\
        \filasep
        Aumento de datos & Recorte aleatorio (70--100\,\% del área), volteo horizontal, RandAugment ($N = 2$) y borrado aleatorio ($p = 0.25$) \\
        \filasep
        Inferencia & \gls{tta} con la imagen y su reflejo horizontal \\
        \filasep
        Precisión numérica & \emph{bfloat16} (bf16) \\
        \bottomrule
    \end{tabularx}
\end{table}

\subsection{Procedimiento experimental}
\label{sec:met_supervisados_ensamble}

Además de los valores comunes, tres hiperparámetros se eligen por búsqueda: la tasa de aprendizaje, el decaimiento de
pesos y la magnitud de RandAugment (Tabla~\ref{tab:met_barrido}). Las cinco tasas se cruzan con los tres
decaimientos, y cada magnitud se asigna a cinco de esas 15 combinaciones. Cada combinación se entrena con y sin pesos
por clase en la función de pérdida (Sección~\ref{sec:marco_desbalance}), lo que suma 30 configuraciones por modelo,
con una ejecución cada una. La búsqueda se hizo sobre Swin-Tiny y ConvNeXt-Tiny, y en ambos ganó, con el criterio de la
Sección~\ref{sec:met_seleccion}, una configuración sin pesos por clase. DenseNet121 y ViT-B/16, en cambio, usan la
configuración ganadora del modelo de su misma familia: la de ConvNeXt-Tiny en el caso de DenseNet121 y la de Swin-Tiny
en el de ViT-B/16. Así los cuatro se comparan con el mismo procedimiento, aunque si DenseNet121 o ViT-B/16 rinden
menos, puede deberse a que la configuración heredada no es la mejor para ellos.

\begin{table}[H]
    \centering
    \caption{Hiperparámetros explorados en la búsqueda y valores elegidos. DenseNet121 usa los de ConvNeXt-Tiny, y ViT-B/16, los de Swin-Tiny.}
    \label{tab:met_barrido}
    \small
    \begin{tabularx}{\textwidth}{@{}>{\raggedright\arraybackslash}p{3.4cm}>{\raggedright\arraybackslash}X>{\raggedright\arraybackslash}p{2.0cm}>{\raggedright\arraybackslash}p{2.2cm}@{}}
        \toprule
        \textbf{Parámetro} & \textbf{Valores explorados} & \textbf{Swin-Tiny} & \textbf{ConvNeXt-Tiny} \\
        \midrule
        Tasa de aprendizaje & $3.0 \times 10^{-5}$; $7.5 \times 10^{-5}$; $1.5 \times 10^{-4}$; $2.5 \times 10^{-4}$; $5.0 \times 10^{-4}$ & $7.5 \times 10^{-5}$ & $3.0 \times 10^{-5}$ \\
        \filasep
        Decaimiento de pesos & 0.01; 0.05; 0.10 & 0.10 & 0.05 \\
        \filasep
        Magnitud de RandAugment ($M$) & 6; 9; 12, cada valor en cinco de las 15 combinaciones & 6 & 9 \\
        \filasep
        Pesos por clase & Inversos a la frecuencia; ninguno & Ninguno & Ninguno \\
        \bottomrule
    \end{tabularx}
\end{table}

Una vez elegida la configuración, cada modelo se entrena tres veces, con distinta inicialización aleatoria, porque las
métricas de las clases minoritarias dependen de pocas imágenes y una sola ejecución no muestra cuánto varía el
resultado.

Los pesos por clase no son la única forma de abordar el desbalance (Sección~\ref{sec:marco_desbalance}), y por eso se
comparan aparte cuatro funciones de pérdida sobre Swin-Tiny, ConvNeXt-Tiny y DenseNet121:
\begin{itemize}
    \item \textbf{Entropía cruzada sin pesos}, la de la configuración común.
    \item \textbf{Entropía cruzada ponderada}, con pesos inversos a la frecuencia de cada clase.
    \item \textbf{Pérdida focal} con $\gamma = 2$~\cite{lin2017focal}, combinada con esos mismos pesos.
    \item \textbf{\emph{Balanced softmax}}~\cite{ren2020balanced}, sin pesos.
\end{itemize}
Para que la función de pérdida sea lo único que cambia, esta comparación fija el resto de los hiperparámetros
(decaimiento de pesos 0.05, $M = 9$ y una tasa de $10^{-4}$, o de $7.91 \times 10^{-5}$ en DenseNet121) y entrena
cada combinación una sola vez.

Por último, los dos modelos de la búsqueda, Swin-Tiny y ConvNeXt-Tiny, también se combinan por voto suave
(\emph{soft voting}; Sección~\ref{sec:marco_senales_confianza}), que promedia sus probabilidades. Solo se combinan esos
dos porque son los que eligió la validación, y elegir otra combinación exigiría consultar la prueba. Para la detección
de melanoma, el puntaje es la probabilidad de melanoma de cada modelo o del ensamble, con el umbral que fija la
Sección~\ref{sec:met_seleccion}.

Estos clasificadores devuelven una probabilidad por clase, pero no generan texto. Para saber si un \gls{vlm}, que sí lo
genera, clasifica tan bien como ellos, se los compara con las mismas métricas (Sección~\ref{sec:met_clasif_vlm}).
